\subsection{The two-primary theta locus is rational over a small field}
\label{sec:two-primary-theta-field}

\paragraph{Finding A35 (a finite constant-field bound).}
Let $C/\mathbf F_q$ be a smooth proper geometrically connected
genus-two curve in characteristic two. Put $X=C^{(1)}$, let $B$ be
the locally exact differential line bundle on $X$, and write
\[
 J=J_X,\qquad
 \Theta=\{L\in J:h^0(X,B\otimes L)>0\}
       =\{\mathcal O_X(R)\otimes B^{-1}:R\in X\}.
 \tag{T2F.1}
\]
Thus $B^{\otimes2}\simeq\omega_X$, with the same Frobenius-twist
normalization as A11 and A24. Let $f\in\{0,1,2\}$ be the
geometric $2$-rank, let $F$ denote $q$-power Frobenius, and let
$d$ be the order of its action on the finite geometric etale group
$J[4](\overline{\mathbf F}_q)\simeq(\mathbf Z/4)^f$; take $d=1$
if $f=0$. Then
\[
 \Theta(\overline{\mathbf F}_q)\cap J[2^\infty]
              \subset J(\mathbf F_{q^d}),
 \qquad
 \begin{cases}
 d=1,&f=0,\\
 d\in\{1,2\},&f=1,\\
 d\in\{1,2,3,4,6\},&f=2.
 \end{cases}
 \tag{T2F.2}
\]
Here $J[2^\infty]$ denotes geometric points of two-power order,
not the full nonreduced group scheme. In particular the rank-one
locus is rational over $\mathbf F_{q^2}$. In every rank it is
rational over $\mathbf F_{q^{12}}$; for an ordinary curve the
sharper field in (T2F.2) has degree at most six, but need not be a
subfield of $\mathbf F_{q^6}$.

The proof uses the two-torsion theta intersection and an elementary
matrix squaring argument. The general bound of two on a nonzero
theta translate is the genus-two Riemann--Roch argument already
used in A19 and A24. Recovered successor A.95 identifies the
rank-one theta characteristic explicitly in its $(3,1)$
Artin--Schreier model \cite[A.95]{ASrun}; no such model or rational
Weierstrass point is assumed here.

\begin{proof}
\emph{1. Nonzero two-torsion translates meet theta only in
two-torsion points.}
Work over $\overline{\mathbf F}_q$. The separable hyperelliptic
double map of $X$ has exactly $f+1$ distinct branch points, by
the Artin--Schreier $2$-rank formula
\cite[Lemma~2.6]{PriesZhu2012}. Write their points on $X$ as
$W_0,\ldots,W_f$. Each fibre divisor at such a point is $2W_i$,
so $2W_i\sim K_X$.

Every difference $[W_i-W_j]$ with $i\ne j$ is a nonzero
two-torsion point. It is nonzero because the Abel map of a curve
of positive genus is injective on geometric points. If $f=1$
this supplies the unique nonzero point of $J[2]$. If $f=2$ the
three unordered pairs supply three distinct points: equality
between two pairs sharing one index would make the other two
Weierstrass points equal in the Abel map, since signs do not
matter for two-torsion. These are all the nonzero points of
$J[2]\simeq(\mathbf Z/2)^f$.

Fix $T=[W_i-W_j]\ne0$. If
$z=\mathcal O_X(R)\otimes B^{-1}\in\Theta$ and
$z+T=\mathcal O_X(S)\otimes B^{-1}\in\Theta$, then
\[
                        R+W_i\sim S+W_j.
 \tag{T2F.3}
\]
If these two effective divisors are equal, their distinct fixed
points $W_i,W_j$ force $R=W_j$ and $S=W_i$. If they are
different, their degree-two class has at least two sections.
Riemann--Roch in genus two then forces that class to be canonical.
The canonical pencil is the hyperelliptic pencil; its fibre
through $W_i$ is $2W_i$. Thus $R=W_i$ and $S=W_j$.
Conversely both displayed possibilities do occur. Consequently
\[
 \Theta\cap(\Theta+T)
   =\{\mathcal O_X(W_i)\otimes B^{-1},
       \mathcal O_X(W_j)\otimes B^{-1}\}
                  \subset J[2]
 \quad\hbox{on geometric points}.
 \tag{T2F.4}
\]
The last inclusion follows from $2W_i\sim K_X\sim2B$.
This formulation uses the actual Raynaud theta divisor. On an
ordinary curve $B$ is non-effective, and replacing $B$ by a
Weierstrass point without accounting for the translation would
give a different theta divisor.

\emph{2. Frobenius cannot move a two-primary theta point.}
For $f=0$ the only geometric two-primary point is zero, so the
claim is immediate. Suppose $f>0$ and identify the geometric
etale Tate module with $\mathbf Z_2^f$. Its associated torsion
group is $(\mathbf Q_2/\mathbf Z_2)^f$, with the compatible
Frobenius action. By the definition of $d$, the matrix of $F^d$
has the form
\[
                        U=I+A,\qquad A\in4M_f(\mathbf Z_2).
 \tag{T2F.5}
\]
Suppose $z\in\Theta$ has exact order $2^a$ and is not fixed by
$U$. Write $z=2^{-a}v$ modulo $\mathbf Z_2^f$, with $v$ a
primitive integral vector. The minimum valuation of the entries
of $Av$ is a finite integer $t<a$, and $t\ge2$.

Set $A_j=U^{2^j}-I$. Induction using
\[
                         A_{j+1}=2A_j+A_j^2
 \tag{T2F.6}
\]
shows both
\[
 A_j\in2^{j+2}M_f(\mathbf Z_2),\qquad
                         v_2(A_jv)=t+j.
 \tag{T2F.7}
\]
Indeed the first assertion is immediate at $j=0$ and is preserved
by (T2F.6). For the second, $2A_jv$ has valuation $t+j+1$, while
$A_j^2v$ has valuation at least
$(j+2)+(t+j)=t+2j+2\ge t+j+2$. Hence their leading terms
cannot cancel, proving the asserted equality at $j+1$.

Choose $j=a-t-1\ge0$. Equation (T2F.7) says that
\[
                    T:=F^{d2^j}z-z
 \tag{T2F.8}
\]
is a nonzero point of exact order two. The line bundle $B$ and
its theta divisor are defined over $\mathbf F_q$, so both $z$
and $z+T$ lie on $\Theta$. By (T2F.4), $z\in J[2]$. But $U$
acts identically on $J[4]$, hence on $J[2]$, contradicting the
choice of $z$. Every two-primary theta point is therefore fixed
by $F^d$, which proves the inclusion in (T2F.2).

\emph{3. The small extension degrees.}
For rank one the automorphism group of the geometric four-torsion
is $(\mathbf Z/4)^\times$, of order two. For rank two, reduction
\[
                  \mathrm{GL}_2(\mathbf Z/4)
                          \longrightarrow\mathrm{GL}_2(\mathbf F_2)
\]
has a kernel of exponent two: $(I+2H)^2=I$ modulo four.
The target has element orders one, two, and three. Therefore
every element upstairs has order dividing two, four, or six.
This gives exactly the asserted list of possible $d$ and proves
all the field bounds.
\end{proof}

\paragraph{The precise finite residual set.}
For $R\in X(\overline{\mathbf F}_q)$ put
$z_R=\mathcal O_X(R)\otimes B^{-1}$. Its canonical class is
\[
                    \kappa_R=[K_X-2R]=-2z_R.
 \tag{T2F.9}
\]
Thus $\kappa_R$ is two-primary if and only if $z_R$ is
two-primary. By (T2F.2) and injectivity of the Abel map defined
by $B$, every such $R$ belongs to $X(\mathbf F_{q^d})$.
An exact enumeration over that one finite field therefore
captures this entire geometric residual set, even when arbitrary
constant extensions are allowed afterwards. The corresponding
statement for the original curve follows by the Frobenius twist.
Also $[R-\iota R]=2z_R$, so the same bound applies to those
hyperelliptic-pair differences that are two-primary.

This is a finite candidate statement, not a separation theorem.
It gives no bound of this kind for arbitrary two-primary
differences $[R-S]$, and it does not treat canonical classes
with a nonzero odd-primary part supported on primes at which
the order of two is odd. Those cases remain distinct from the
purely two-primary residual set considered here. No GRH or
prime-distribution assertion enters the proof.

\paragraph{The ordinary real-multiplication targets need only degree six.}
Suppose in addition that $C$ is ordinary and its Jacobian has an
action of $\mathcal O_5=\mathbf Z[(1+\sqrt5)/2]$ defined over
$\mathbf F_q$, as do the ordinary reflection targets in A33.
The action carries over to $J_X$ by Frobenius twist. Then
\[
 \Theta(\overline{\mathbf F}_q)\cap J[2^\infty]
                     \subset J(\mathbf F_{q^6}),
 \qquad
 \{R:\kappa_R\text{ is two-primary}\}
                     \subset X(\mathbf F_{q^6}).
 \tag{T2F.10}
\]
Here is the additional argument. The polynomial $T^2+T-1$ is
irreducible modulo two, so
$\mathcal O=\mathcal O_5\otimes\mathbf Z_2$ is the ring of
integers of the unramified quadratic extension of $\mathbf Q_2$.
The geometric etale Tate module has rank two over $\mathbf Z_2$.
It is a finite torsion-free $\mathcal O$-module: a nonzero
$\mathcal O$-annihilator would, on taking its norm, give a
nonzero $\mathbf Z_2$-annihilator. Since $\mathcal O$ is a
discrete valuation ring, the module is free, necessarily of
rank one over $\mathcal O$.

Frobenius commutes with this action because the endomorphisms are
defined over $\mathbf F_q$. It therefore acts by a unit of
$\mathcal O$. Reduction of $(\mathcal O/4)^\times$ modulo two
has quotient $\mathbf F_4^\times$, of order three, and kernel
$1+2\mathcal O$ modulo four, of exponent two. Thus
$(\mathcal O/4)^\times$ has exponent six, and $F^6$ acts
identically on $J[4]$. Applying the proof of (T2F.2) with this
power gives (T2F.10). This confines the purely two-primary case
on the reduced targets to one explicit finite field; it still
does not force a canonical class to have a good odd-primary
component or separate any of these remaining points.
